% MAS1004 Assignment 1: long report.
%
% You fill this in with your agent, in your repository. AGENTS.md tells the
% agent how. Red text in the PDF is what still has to be replaced.
%
% - Pictures go in figures/. Copy them from results/, which is not in git.
% - Outputs go in outputs/. run.py saves what it prints in
%   results/<tag>_output.txt, and PLAN.md saves the other outputs in results/
%   with `| tee results/...`. Copy them here. \outputfile{outputs/check.txt}
%   prints a file exactly as it is. Every number in this report has to come
%   from such an output.
%
% To make the PDF: compress this folder into a zip, and on Overleaf choose
% New Project, Upload Project. Set Menu, Compiler to XeLaTeX, and Recompile.

\documentclass[11pt]{article}
\usepackage{mas1004}

\reporttitle{Assignment 1: Long report}
\studentname{김서인}
\studentid{20251121}

\begin{document}
\makeheader

% ---------------------------------------------------------------------------
\section*{Problem 1. Categories and images}

\subsection*{Why these categories}
I run the planning/design side of an IT club, so I spend a lot of time
hunting for design references. Pinterest only groups images by mood, not by
art style, so telling a pixel-art reference from a hand-drawn one from a
vector one was something I always had to judge myself, one image at a time.
I chose these three categories so that an automatic style classifier could do
that judgement for me, and help me decide which references are actually
useful before I open them.

\subsection*{What I expected the model to use}
Before training, I expected pixel art to be easy: the visible pixel grid
(blocky, stair-stepped edges) is a texture that has nothing to do with the
other two. I expected hand-drawn doodles and vector illustration to be the
hard pair, told apart by how regular the outline is: smooth, uniform-width
vector curves versus the shaky, uneven line weight of a hand-drawn line.

\subsection*{Images per category}
150 images were downloaded for each of the three categories before any
cleaning:
\begin{itemize}
  \item \texttt{ink\_drawing\_doodle\_sticker} (hand-drawn): 150
  \item \texttt{cute\_pixel\_art\_sticker\_design} (pixel art): 150
  \item \texttt{kawaii\_sticker\_vector\_sheet} (vector illustration): 150
\end{itemize}

% ---------------------------------------------------------------------------
\section*{Problem 2. Preparing an image}

\sidebyside{figures/prepare_original.jpg}{the image as downloaded (1920\,$\times$\,592)}
           {figures/prepare_result.png}{what prepare\_image made of it (224\,$\times$\,224)}

This doodle sheet is 1920 by 592 pixels, more than three times wider than it
is tall. Resizing the shorter side to 256 and cutting out the centre 224 by
224 square throws away almost everything on the left and right: a rabbit face
in the top-right corner and a bear face in the bottom-left corner are both
completely cut out of what the model actually sees, even though they are
clearly part of the sheet.

% ---------------------------------------------------------------------------
\section*{Problem 3. First training runs}

\begin{center}\footnotesize
\begin{tabular}{llllrrrr}
\toprule
tag & start & trained & epochs & lr & trainable parameters & train acc. & test acc. \\
\midrule
run     & ImageNet & all layers  & 10 & 0.0001 & 11,178,051 & 99.4\%  & 93.3\% \\
scratch & random   & all layers  & 10 & 0.0001 & 11,178,051 & 100.0\% & 74.4\% \\
frozen  & ImageNet & last layer  & 10 & 0.001  & 1,539      & 91.7\%  & 90.0\% \\
short   & ImageNet & all layers  & 5  & 0.0001 & 11,178,051 & 100.0\% & 94.4\% \\
\bottomrule
\end{tabular}
\end{center}

\picturefile{figures/run_curves.png}

The training loss drops sharply for the first 2--3 epochs and then flattens
out near zero. The test loss follows it down for about 3 epochs, bottoms out
around epoch 3 (0.151), and then slowly climbs back up to 0.259 by epoch 10
even though the training loss keeps falling -- a textbook overfitting curve.

Train accuracy (99.4\%) is 6.1 points higher than test accuracy (93.3\%) on
the \texttt{run} setting. The model has nearly memorised the training images
by the end of training, so the gap is the model fitting quirks of the
training images that do not generalise, not real extra ability.

\begin{center}\picturefile[0.5\linewidth]{figures/run_confusion.png}\end{center}

The two categories it mixes up most are hand-drawn doodle and vector
illustration: 5 of the 30 hand-drawn test images were called vector, while
pixel art was never confused with either of the other two (30/30 correct).
This did not surprise me -- it is exactly the pair I expected to be hardest,
since both styles can have fairly clean line art once a hand-drawn sketch
happens to be inked neatly.

Starting from ImageNet reached 93.3\% test accuracy; starting from random
numbers reached 74.4\% -- 19 points apart, with everything else (images,
code, epochs) identical. I think this is because the ImageNet weights
already encode general-purpose edge, shape and texture detectors from 1.2
million photographs, so the network only has to learn how those relate to
our three sticker styles, instead of learning to see at all from scratch.

Training only the last layer changes 1,539 parameters, about 0.0138\% of the
11,178,051 total. It still reached 90.0\% test accuracy, only 3.3 points
behind training every parameter (93.3\%) -- surprisingly close, given how
tiny a fraction of the network it is allowed to touch.

\subsection*{Output of the runs}
\outputfile{outputs/run_output.txt}
\outputfile{outputs/scratch_output.txt}
\outputfile{outputs/frozen_output.txt}
\outputfile{outputs/short_output.txt}

% ---------------------------------------------------------------------------
\section*{Problem 4. Cleaning}

\subsection*{My rule}
Common rule, applied to all three categories: remove any image that is a
real-life photo of a sticker (rather than the illustration itself), a
screenshot of design software, or an ad/product-listing graphic where banner
text, a price tag or a watermark takes up a large share of the picture.

Per-category rule, on top of the common rule: for \emph{pixel art}, remove
anything without a visible pixel grid (smooth curves or gradients instead).
For \emph{hand-drawn}, remove anything with a uniform line weight and flat
colour fill and no hand-drawn irregularity. For \emph{vector}, remove
anything that instead looks like pixel art (visible grid) or like a hand-drawn
doodle (uneven, sketchy line).

\subsection*{Three kinds of junk}

\begin{center}
\picturefile[0.3\linewidth]{figures/junk_real_photo.jpg}\hfill
\picturefile[0.3\linewidth]{figures/junk_ad_graphic.jpg}\hfill
\picturefile[0.3\linewidth]{figures/junk_near_duplicate.jpg}
\end{center}

Left to right: a real photo of a hand holding a physical sticker sheet
(common-rule violation, not the illustration itself); a product-listing
graphic with a ruler, size labels and a logo wrapped around the actual
stickers (common-rule violation, mostly ad text); and one half of a
near-duplicate pair, the same sticker sheet photographed/re-saved twice.

\subsection*{What I removed}
\outputfile{outputs/clean_count.txt}

\texttt{clean.py suspects} found 67 pairs of near copies (similarity above
0.95). Several images were near-copies of more than one other image at once,
so the 67 pairs collapse into 44 groups once images that are copies of each
other are merged together; keeping only one image per group removed 55
files.

Looking at the top 20 suspects (the 20 images in each class that look least
like the rest of the class) by my rule: I removed 6/20 in pixel art, 5/20 in
hand-drawn, and 7/20 in vector. Several of these were junk I had not
separately flagged when I first looked through the sheets in Step 3 (small
real-photo and ad-banner images are easy to miss among 150 thumbnails), so
the suspect list genuinely caught things I had missed. On the other hand,
most of what I actually removed -- the 55 near-duplicate images -- came from
the separate \texttt{copies.png} list, not from the "looks least like its own
class" suspect ranking, so the suspect ranking alone would have missed the
single largest source of junk in this dataset.

\subsection*{Before and after cleaning}
\outputfile{outputs/compare.txt}

On the (different, smaller) downloaded test set, test accuracy went down
slightly after cleaning: 93.3\% (\texttt{run}) to 90.7\% (\texttt{clean}).
That is not the fair comparison, though, since cleaning changes which images
are in the test set. On the new-source images, which did not change, accuracy
stayed exactly the same before and after cleaning: 47.4\% (9/19) both times.
So cleaning mainly removed noise and near-duplicates that were making the
model's own, now-smaller test set easier -- it neither helped nor hurt how
well the model generalises to genuinely new images.

% ---------------------------------------------------------------------------
\section*{Problem 5. Images from a new source}

My new images came from Pinterest: I searched and saved sticker-sheet images
by hand, picking specific pins rather than the first generic result, since
Pinterest content is also reachable by the DuckDuckGo/Bing search that
\texttt{collect.py} uses. I then ran \texttt{clean.py overlap} to check.

\outputfile{outputs/overlap.txt}

None of the 19 new images is a near copy of any of the 830 downloaded
images; the closest match anywhere is 0.93 similarity, below the 0.95
near-copy threshold.

Accuracy on the downloaded test set (clean model): 90.7\%. Accuracy on the 19
new images: 47.4\% (9/19) -- much worse, as the assignment predicts.

\textbf{Five confident mistakes:}

\sidebyside{figures/mistake1_pixel_as_ink.png}{really pixel art; said hand-drawn doodle (100\%)}
           {figures/mistake2_ink_as_vector.png}{really hand-drawn; said vector (100\%)}

\sidebyside{figures/mistake3_ink_as_vector.png}{really hand-drawn; said vector (100\%)}
           {figures/mistake4_ink_as_vector.png}{really hand-drawn; said vector (99\%)}

\begin{center}\picturefile[0.4\linewidth]{figures/mistake5_pixel_as_vector.png}\end{center}
really pixel art; said vector (95\%)

For the pixel-art image called hand-drawn (top left): it is monochrome,
black-outline-only pixel icons with no colour fill at all, so it visually
reads as a black-and-white line drawing, and I think the model used the lack
of colour rather than noticing the pixel grid. For the three hand-drawn
images called vector: all three have thin, faint outlines (one is almost
pure colour-pencil shading with barely a line at all) -- exactly the kind of
image Problem 7 tests below. For the pixel-art sheet called vector: it uses
soft pastel colours, a smooth gradient background and a clean white sticker
border, which is closer to the soft, clean presentation of our vector
stickers than to the flat, saturated colours of most of the pixel-art
training images.

What is different about my new images: they are full product photographs or
Pinterest pin images rather than isolated sticker cut-outs on a plain
background -- they include borders, packaging text, soft gradient
backgrounds, and sometimes a hand holding the sticker, none of which appear
much in the search-engine downloads.

% ---------------------------------------------------------------------------
\section*{Problem 6. The demo}

Address: \url{https://jzadoo.github.io/mas1004-assignment1/}

When I fed the model a 300\,$\times$\,300 image of pure random colour noise
(nothing like any of the three styles), it still had to pick one of the three
answers: Pixel Art 51\%, Hand-drawn Doodle 45\%, Vector Illustration 4\%. The
model has no "none of the above" option, so on something that is none of its
categories it just guesses, roughly split between the two "busy, high-detail"
categories and almost never the "smooth, flat" vector category.

% ---------------------------------------------------------------------------
\section*{Problem 7. What the model looks at}

\subsection*{Guess 1}

Looking at \texttt{results/clean\_worst.png} and the new-image mistakes above,
hand-drawn doodles kept being called vector, and the ones I checked by eye
all had thin, faint outlines. My guess: the model tells hand-drawn doodles
apart from vector illustration mainly by how thick/visible the outline is,
not by how irregular or "hand-drawn" the line looks -- a thin but shaky line
might still read as vector if the model simply cannot resolve it at
224$\times$224.

I chose to thicken the outlines, because if the model is missing thin lines
rather than reading their shakiness, making them thicker (without making them
any smoother or more regular) should fix the mistake even though the guess
about shakiness is unrelated to thickness.

My request to the agent, word for word:
\begin{quote}
"선을 두껍게 하는 거" (make the outlines thicker) -- applied as a MinFilter
(size 5) on every image in \texttt{data/new\_images/ink\_drawing\_doodle\_sticker},
saved into \texttt{data/changed/ink\_drawing\_doodle\_sticker\_thick\_lines}.
\end{quote}

\sidebyside{figures/guess1_original_001.png}{original: vector (wrong)}
           {figures/guess2_changed_001.png}{thick lines: hand-drawn (correct)}

\sidebyside{figures/guess1_original_004.png}{original: vector (wrong)}
           {figures/guess2_changed_004.png}{thick lines: hand-drawn (correct)}

\sidebyside{figures/guess1_original_006.png}{original: vector (wrong)}
           {figures/guess2_changed_006.png}{thick lines: vector (still wrong)}

Across all 8 new hand-drawn images, accuracy went from 2/8 (25\%) on the
originals to 7/8 (87.5\%) after thickening the outlines -- a large jump.
Image 006 is the one exception: looking at it directly, it is almost pure
colour-pencil shading with barely any ink outline to thicken in the first
place, which fits the guess rather than contradicting it (there was no thin
line there -- there was close to no line at all).

\subsection*{Guess 2}

Before settling on thickening, I first tested blurring the same 8 images,
expecting that removing fine line detail would make hand-drawn doodles look
smoother and more vector-like.

My request to the agent, word for word:
\begin{quote}
"블러일까?" (would blur do it?) -- applied as a Gaussian blur (radius 4px) on
every image in \texttt{data/new\_images/ink\_drawing\_doodle\_sticker}, saved
into \texttt{data/changed/ink\_drawing\_doodle\_sticker\_blurred}.
\end{quote}

\sidebyside{figures/guess1_original_004.png}{original: vector (wrong)}
           {figures/guess1_changed_004.png}{blurred: pixel art (wrong, different way)}

\sidebyside{figures/guess1_original_002.png}{original: vector (wrong)}
           {figures/guess1_changed_002.png}{blurred: hand-drawn (correct)}

Blurring barely changed overall accuracy (2/8 to 3/8), but it changed
\emph{which} mistake the model made: every image that was originally called
vector and stayed wrong was now called pixel art instead, not vector. My
reading is that blur does not make a doodle's lines look smooth and
vector-like -- it turns them into soft, grainy colour blobs, which the model
associates more with pixel art's busy small shapes than with vector's large
flat regions. This guess, on its own, does not explain the mistakes as
cleanly as Guess 1 does, which is why I moved on to thickening the lines
instead.

\subsection*{What I now think the model uses}
I am fairly confident, based on the 25\%\,$\to$\,87.5\% jump from thickening
alone, that outline strength/visibility is the main thing separating
hand-drawn doodle from vector illustration for this model -- not line
regularity, which was my original guess before training. Pixel art is told
apart from the other two by its grid texture most of the time, but the one
pixel-art mistake above shows that a monochrome, outline-only pixel sheet can
still be read as hand-drawn when there is no colour fill to anchor it as
pixel art.

% ---------------------------------------------------------------------------
\section*{Appendix. Output of check.py}
\outputfile{outputs/check.txt}

\end{document}
